\section{Discussion}

\subsection{Deployment Implications}

Our evaluation suggests that behavioral attestation is practical for the class
of HPC agent workloads targeted by \projectname, but its value depends on deployment
discipline rather than on instrumentation alone. The most important lesson is
that agent identity is not a sufficient security primitive for autonomous
scientific workflows. Once an agent is granted legitimate filesystem, network,
and tool access, the relevant question is whether its \emph{observed behavior}
continues to match its authorized task. \projectname~therefore fits naturally into HPC
sites that already operate with strong identity and scheduler control, but need
a runtime mechanism for detecting abuse of legitimate authority.

The real-path results in Section~\ref{results} show that this model can be realized with an
\textit{ebpf}-based probe, a verifier-driven policy pipeline, and scheduler-mediated
containment without imposing large steady-state cost. The 1\,s configuration
provides sub-second median detection latency with low measured verifier CPU
overhead, which makes it a practical default for interactive and batch-oriented
agent workflows. More importantly, the experiments show that cluster context is
necessary: the coordinated attack is only visible when the verifier correlates
behavior across agents rather than reasoning about each session in isolation.

In practice, this means \projectname~should be deployed as part of the control plane,
not as an application library. The constraint manager, verifier, and containment
logic need access to job metadata, scheduler actions, and cluster-wide event
streams. The design is therefore most useful in settings where the site can
control both the job launch path and the enforcement path, such as Slurm-based
research clusters, AI-enabled scientific gateways, or managed institutional
agent platforms.

\subsection{Limitations}

The primary limitation of behavioral attestation is that its guarantee is only
as strong as the behavioral constraint profile. If a profile omits an important
resource dimension or encodes an overly permissive boundary, then actions in
that gap are not violations. This is not a flaw in the verifier; it is a policy
specification problem. Templates and inference can reduce operator burden, but
high-assurance deployments still require human review of the resulting
constraints.

A second limitation is response granularity. \projectname~does not prevent the initial
attack step; it detects and contains the resulting policy violation. The size of
that response window is bounded by the attestation interval and the containment
path. Our evaluation shows that the trade-off is favorable at a 1\,s interval,
but the window is not zero, and highly latency-sensitive deployments may need
shorter intervals or stronger inline enforcement for selected resources.

A third limitation is the trust model. The current design assumes a trustworthy
kernel, a trustworthy node-level collector, and an uncompromised scheduler
control plane. \textit{ebpf}-based monitoring raises the bar relative to user-space
instrumentation, but it does not solve the problem of a malicious kernel or a
fully privileged attacker on the node. Similarly, the centralized verifier used
in our prototype simplifies cluster-wide correlation, but it is not yet the
final scaling architecture for very large systems with thousands of nodes and
high event rates.

Finally, our evaluation intentionally separates real-path measurements from the
larger controlled studies. This is the right choice for methodological clarity,
but it means some claims in Section~\ref{results} are deployment measurements, whereas others
are controlled security results. We view that split as a strength rather than a
weakness, but it also defines the boundary of what the current prototype has
physically exercised.

\subsection{Future Directions}

Several extensions would strengthen the system. First, constraint generation can
be improved with richer policy templates, stronger provenance for inferred
profiles, and review workflows that make it easier for users and site operators
to validate what an agent is actually allowed to do. Second, the verifier can be
extended with more scalable data structures and distributed aggregation for
cluster-wide correlation. Third, containment can be made more adaptive by tying
verdicts to job QoS controls, storage isolation, and site-specific incident
response policies.

More broadly, behavioral attestation appears useful beyond the specific threat
model studied here. The same core idea, declare what an autonomous actor is
allowed to do, continuously verify observed behavior, and contain violations at
runtime, applies to other multi-tenant AI deployments, including cloud-hosted
agents, scientific gateways, and tool-using assistants operating over shared
enterprise data. HPC is a compelling first target because it combines valuable
data, strong scheduler control, and increasingly autonomous agents, but the
underlying security primitive is more general.
